\documentclass[11pt,a4paper]{article}
\usepackage[margin=1in]{geometry}
\usepackage{hyperref}
\usepackage{enumitem}
\usepackage{graphicx}
\usepackage{xcolor}

% -----------------------------------------------------
% bvraghav-tiet-ucs503p-article.sty
% -----------------------------------------------------

% Variable: Lab Instructor
\makeatletter \newcommand{\BvrSetLabInstructor}[1]{%
\gdef\Bvr@LabInstructor{#1}}
% default
\newcommand{\Bvr@LabInstructor}{Dr. Paramveer Kaur}
% public accessor
\newcommand{\BvrLabInstructorValue}{\Bvr@LabInstructor}
\makeatother

% Add subtitle
\usepackage{titling}
% -- Set title bold
\pretitle{\begin{center}\bfseries\fontsize{18bp}{18bp}\selectfont}
\makeatletter
\newcommand{\BvrSubtitle}[1]{%
\posttitle{%
\par\end{center}
\begin{center}\large#1\end{center}
\vskip 0.5em}%
}
\makeatother

% Hints in the section title(s)
\newcommand{\BvrHint}[1]{%
{\normalfont\normalsize\itshape\hfill #1}}

% -----------------------------------------------------

\title{ArogyaMitra: An AI-Assisted Real-Time \ Telemedicine Consultation Platform}

\BvrSetLabInstructor{Dr. Paramveer Kaur}

\BvrSubtitle{%
Project Proposal for UCS503P  \\
Submitted to: \BvrLabInstructorValue \\ \bfseries
Thapar Institute of Engineering and Technology%
}

\author{
Animesh Sudhanshu -
{ \texttt{1024170375}}
\and
Ashish Bhagat -
{\texttt{1024170372}}
\and
Devansh Rathaur - 
{\texttt{1024170380}}
}

\date{\today}

\begin{document}
\maketitle

\section{Higher-order goal%
\BvrHint{\ldots secondary framing}}
This project aims to improve access to primary
healthcare by reducing the friction between a patient
needing medical advice and a doctor available to give
it. While the broader motivation is healthcare equity
India's government telemedicine platform), the
immediate engineering objective is to translate
\emph{low-friction, real-time doctor-patient
consultation} into a measurable, deployable software
solution built and iterated on by a small student team.

\section{Time-to-value %
\BvrHint{\ldots primary heuristic}}
\begin{itemize}[leftmargin=0.7cm]
\item We will deliver meaningful  core consultation workflow and validating it within a short iteration window.
\item We will prioritize fast feedback over perfection by using weekly iterations (and targeting CI-backed automation early) to reduce release friction.
\item Success is defined by what can be delivered and measured in the first iteration (a working doctor--patient consultation loop), then improved based on user feedback.
\end{itemize}

\section{Problem Statement}
Patients, particularly those in smaller towns with
limited specialist availability, face delays and
friction when seeking timely medical advice for
non-emergency conditions. Common pain points include:
\begin{itemize}[leftmargin=0.7cm]
\item \textbf{Access friction:} booking an in-person consultation requires travel and waiting, even for simple queries.
\item \textbf{Fragmented communication:} follow-up questions, prescriptions, and advice are often exchanged informally (calls, messaging apps) with no persistent record.
\item \textbf{No triage before consultation:} patients cannot easily get a first-pass sense of urgency or relevant information before a doctor is available.
\item \textbf{Scheduling opacity:} patients cannot easily see doctor availability or track the status of a booked consultation.
\end{itemize}

\textbf{Why this matters now (contextual relevance):}
Real-time web technologies (WebRTC, WebSockets) and
affordable managed backends have matured enough that a
small team can build a credible consultation platform
in an academic timeframe, while conversational AI can
plausibly assist (not replace) the pre-consultation
triage step.

\section{Proposed Solution %
\BvrHint{\ldots Software Proposition}}
\subsection{Overview}
Build a web-based ``ArogyaMitra'' real-time
telemedicine platform
\begin{itemize}[leftmargin=0.7cm]
\item \textbf{Patient portal:} registration, profile management, doctor discovery, and appointment booking.
\item \textbf{Doctor dashboard:} availability management, appointment queue, and consultation history.
\item \textbf{Real-time consultation:} secure real-time chat and video call between doctor and patient.
\item \textbf{AI pre-consultation assistant:} a conversation-aware assistant that helps patients describe symptoms and prepare context before a doctor joins, with a safety layer that always defers emergency-flagged symptoms to immediate care.
\item \textbf{Notifications:} real-time alerts when a consultation is about to begin or a doctor becomes available.
\end{itemize}

\subsection{Core workflow}
\begin{enumerate}[leftmargin=0.7cm]
\item Patient registers/logs in and books a consultation slot with an available doctor.
\item Patient optionally interacts with the AI assistant to summarize symptoms ahead of time.
\item Both parties receive a real-time notification when the consultation is about to start.
\item Doctor and patient connect via real-time chat / video call; doctor records notes and advice.
\item Patient can review consultation history and past notes from their profile.
\end{enumerate}

\subsection{Operational constraints for an educational setting}
\begin{itemize}[leftmargin=0.7cm]
\item Keep the solution usable on low-to-mid range devices via responsive web design.
\item Avoid dependence on novel research algorithms; focus on workflow correctness, real-time reliability, and safe AI-assistance boundaries.
\item Design for deployability using free-tier or low-cost managed hosting suitable for a pilot (e.g., Render, Netlify, MongoDB Atlas).
\end{itemize}

\section{Solution Approach %
\BvrHint{\ldots Engineering focus}}
This is an engineering course project, so we focus on
integration, workflow correctness, and safe deployment
of existing well-understood technologies rather than
novel algorithm research.

\subsection{AI assistant scope %
\BvrHint{\ldots non-research}}
The AI pre-consultation assistant is an integration
task, not a research contribution:
\begin{itemize}[leftmargin=0.7cm]
\item Use an existing LLM API (e.g., Gemini) for conversational symptom intake; no in-house model training.
\item Add a lightweight safety/triage layer that detects emergency-flagged keywords (e.g., chest pain, stroke signs, self-harm) and always routes to an ``seek immediate care'' response rather than open-ended advice.
\item The assistant only prepares context for the doctor; it never issues a diagnosis or prescription.
\end{itemize}

\subsection{Web architecture %
\BvrHint{\ldots web-based solution}}
A typical decoupled web architecture:
\begin{itemize}[leftmargin=0.7cm]
\item \textbf{Frontend:} Next.js with a component library (Shadcn) for patient/doctor portals and the consultation UI.
\item \textbf{Backend API:} Node.js/Express for authentication, appointment logic, and consultation record management.
\item \textbf{Database:} MongoDB Atlas for users, appointments, consultation notes, and message history.
\item \textbf{Real-time layer:} Socket.io for chat/notifications; WebRTC for video calls.
\item \textbf{Authentication:} Firebase Authentication for patient and doctor accounts.
\end{itemize}

\subsection{CI/CD and fast delivery enabler}
CI/CD will be used to:
\begin{itemize}[leftmargin=0.7cm]
\item Automatically build and run tests (lint, unit) on every change via GitHub Actions.
\item Produce repeatable deployment artifacts to a staging environment (Render for backend, Netlify for frontend).
\item Enable safe and frequent releases of incremental features.
\end{itemize}

\section{Evaluation Criterion %
\BvrHint{\ldots measurable and attributable}}
We define success using measurable metrics tied to the
core consultation workflow.

\subsection{Primary evaluation metric}
\textbf{Time-to-connect (TTC):} time between a
scheduled consultation's start time and the moment both
doctor and patient are connected in the same
chat/video session.
\begin{itemize}[leftmargin=0.7cm]
\item Target: median TTC $\le$ 2 minutes during pilot testing.
\item Attribution: measured from server-side timestamps (scheduled time vs. session-join event).
\end{itemize}

\subsection{Secondary evaluation metrics}
\begin{itemize}[leftmargin=0.7cm]
\item \textbf{Booking completion rate:} percentage of started bookings that result in a confirmed appointment.
\item \textbf{AI assistant safety accuracy:} percentage of emergency-flagged test cases correctly routed to ``seek immediate care'' during a fixed test suite run.
\item \textbf{User clarity:} patient-reported clarity of the consultation flow using a simple 1--5 feedback question.
\item \textbf{Reliability:} availability of staging/production for login, booking, and consultation during test windows (e.g., $\ge$ 99% uptime in pilot).
\end{itemize}

\subsection{Pilot validation plan %
\BvrHint{\ldots high-level}}
\begin{itemize}[leftmargin=0.7cm]
\item Recruit 10--15 pilot patients and 3--5 volunteer/simulated doctor accounts.
\item Run a fixed AI-assistant test suite (symptom prompts including emergency cases) before each release to catch safety regressions.
\item Compare TTC and booking completion rate across iterations.
\end{itemize}

\section{Scalability %
\BvrHint{\ldots with foresight}}
The proposition should scale in both theoretical and
practical senses.

\begin{enumerate}[leftmargin=0.7cm]
\item \textbf{Scalable (theoretically):} The system should support growth in number of concurrent consultations and users by:
\begin{itemize}[leftmargin=0.7cm]
\item decoupling frontend and backend (already true of the Next.js/Express split),
\item using stateless REST endpoints with Socket.io namespaces scoped per consultation room,
\item indexing appointment and user queries in MongoDB.
\end{itemize}

\item \textbf{Operationally deployable (practically):} The system should run on common managed hosting:  
\begin{itemize}[leftmargin=0.7cm]  
    \item MongoDB Atlas as a managed database,  
    \item Render/Netlify as platform-hosted deployment targets,  
    \item CI/CD pipeline for staging/production via GitHub Actions.  
\end{itemize}

\end{enumerate}

\section{Engine availability heuristic}
The team focuses on workflow and safety
correctness:
\begin{itemize}[leftmargin=0.7cm]
\item Next.js for the frontend framework.
\item Node.js/Express for REST API backend.
\item MongoDB Atlas as the managed document database.
\item Firebase Authentication for identity.
\item Socket.io for real-time chat/notifications; WebRTC for video.
\item Gemini API / Vercel AI SDK for the conversational assistant.
\item GitHub Actions as the CI runner; Render/Netlify as deployment targets.
\end{itemize}

\section{Project scope and deliverables %
\BvrHint{\ldots iteration-friendly}}
\subsection{Initial deliverable %
\BvrHint{\ldots first iteration}}
\begin{itemize}[leftmargin=0.7cm]
\item Set up local + staging environments.
\item Patient and doctor authentication and profile pages.
\item Appointment booking flow with doctor availability.
\item Real-time chat consultation between doctor and patient.
\item Basic logging and timestamps for TTC measurement.
\item CI: build + lint + backend unit tests.
\item CD to staging (Render + Netlify) with a single pipeline job.
\end{itemize}

\subsection{Subsequent deliverables}
\begin{itemize}[leftmargin=0.7cm]
\item Real-time video call integration (WebRTC) alongside chat.
\item AI pre-consultation assistant with the emergency-triage safety layer.
\item Real-time notifications for upcoming/available consultations.
\item Admin/analytics view of bookings, TTC, and completion rate.
\item Regression test suite for AI assistant safety behaviour.
\end{itemize}

\section{Risks and mitigations}
\begin{itemize}[leftmargin=0.7cm]
\item \textbf{Data privacy / consent:} minimize collected medical data; store only what is required for the consultation record; never persist raw AI assistant conversations beyond what is needed for the doctor's context.
\item \textbf{AI safety failure:} maintain a small, fixed regression test suite of emergency-symptom prompts that must pass before every release; the assistant never issues a diagnosis or prescription.
\item \textbf{Real-time reliability:} use WebRTC/Socket.io connection-state handling and fallbacks (e.g., reconnect logic) so a dropped session does not lose consultation context.
\item \textbf{Staff/doctor adoption:} keep the doctor dashboard minimal (availability toggle, queue, notes) to reduce onboarding friction for simulated/volunteer doctors during the pilot.
\item \textbf{Operational issues in pilot:} use staging early; ensure CI/CD produces repeatable deployments before the pilot window.
\end{itemize}
\section{Summary}
This project proposes a deployable, AI-assisted
real-time telemedicine platform, to
reduce the friction between patients and doctors for
non-emergency consultations. It meets the course
criteria by emphasizing time-to-value through rapid
iterations, implementing CI/CD to
support frequent safe releases, and using measurable
evaluation criteria (especially time-to-connect and AI
assistant safety accuracy). It remains focused on
engineering integration, workflow correctness, and safe
AI-assistance boundaries rather than research-driven
novelty.

\end{document}
